% Chapter 6. Answers to RQs (after the pilot), contributions, doctoral directions (docs/roadmap.md, stage 2; ADR-0051, ADR-0064).
\chapter{Conclusion}\label{ch:concl}

\section{Summary}

The thesis set out to collect a small curated speech dataset of primer-level reading by beginners annotated by
teachers, to evaluate existing speech recognition and analysis models on it, and to use agentic pipeline optimization to build from them a reading and pronunciation error-detection system for an out-of-class practice tool.
The literature review showed that the primer's sound-based, incremental method suits automatic feedback and
provides a ready-made skill map; that nobody had measured the best open model for Qur'anic Arabic on
primer-level reading by second-language beginners or adapted it to them, although small-data methods for this
exist; and that tutors that listen to reading work for children but have not been applied to adults learning to read a foreign script, and that benchmarks do not separate reading errors from pronunciation errors.

The collection tool built in response turns the whole book into 125 lessons and 800 cards traceable to page and
line; grades every voice message at the phoneme level with a full decision trace; explains errors in Russian;
and stores consented, pseudonymised data ready for teacher annotation that separates reading from pronunciation
errors. Agentic pipeline optimization is designed as a controlled experiment: an AI agent searches for and combines
configurations of off-the-shelf models and components --- a calibrated verifier with a ``not sure'' outcome and a controlled
false-rejection rate, LoRA fine-tuning with augmentation from the book's minimal pairs --- only on the validation part, within a fixed budget; every configuration tried is logged, and the best one is compared once with the baseline configurations on a frozen test set annotated by two teachers. All of this was built under mechanical checks
that protect both the code and the experiments.

\section{Answers to the research questions}

\paragraph{RQ-1 (reliability of off-the-shelf models)} \pending{answer after the pilot and double annotation}.
\paragraph{RQ-2 (agentic pipeline optimization)} \pending{answer after the agentic optimization and the single test run}.


\section{Contributions}

\begin{enumerate}
  \item[\textbf{C1.}] A dataset of primer-level reading by beginners learning Arabic as a foreign language, tied to
    the skills of \emph{al-Muallim ath-Thani}, collected with consent, with phoneme-level and error-type annotation
    by teachers, measured reliability, a gold test set and a datasheet.
  \item[\textbf{C2.}] The reliability of open speech recognition and analysis models on beginners, measured
    against teachers.
  \item[\textbf{C3.}] An error-detection system built by agentic pipeline optimization from off-the-shelf models and components, with its gain over baseline configurations, a log of every configuration tried and a leakage-free
    evaluation protocol.

\end{enumerate}

The engineering result is a reproducible ``collect $\to$ annotate $\to$ agentic optimization $\to$ evaluate'' pipeline, a
digital curriculum of the book traceable to an edition and a tutor with an open architecture in which models can
be compared in shadow mode; the methodological result is a documented harness for AI-agent development and ML
experiments.

\section{Author's contribution}\label{sec:concl-author}
% Форма отзыва руководителя (Прил. 14): самостоятельность выполнения. [утвердить] с автором и руководителем.

The author defined the topic and the research questions together with the supervisor; prepared the digital
curriculum of the primer and its skill map; took the architectural decisions recorded in the project's
architecture decision records; wrote the preregistration and the analysis plan; organises the pilot and the
annotation by two teachers; and took every decision on teaching content, texts shown to learners, thresholds,
consent and data. Most of the code was written by an AI coding agent under the author's direction and under
mechanical checks (Section~\ref{sec:impl-harness}); the statement on the use of AI tools gives the details. The
teachers annotate the data, and the agent has no access to test-set labels.

\section{Dissemination}\label{sec:concl-dissemination}
% Форма отзыва (Прил. 14): публикации и доклады по теме ВКР.

\pending{publications and presentations on the topic of the thesis}. The plan is two rounds of expert feedback:
on the framing and the pilot design (December--January) and on the mature evaluation (March--April), at
workshops whose audience works on speech assessment or educational technology, or at a departmental seminar if
no suitable call fits the dates. The dataset, the code and the fine-tuned weights are to be released if ethics
and consent allow (contribution C1).

\section{Future work}

The master's thesis is designed as the first stage of a doctoral project and hands it its tools: the dataset; a
frozen test set that stays a test set, so that future results remain comparable; the skill map; the agentic pipeline optimization and its configuration log; and shadow mode for comparing models on live attempts. The
main doctoral line is \emph{teaching the AI and learning with it}; it splits into the following directions.
\begin{enumerate}
  \item \textbf{Deferred master's questions.} How the AI learns from teacher corrections and the book's
    minimal pairs, whether mastery of each skill can be measured from speech, and how adults progress through
    the book's skills; these questions and their hypotheses were moved out of the master's thesis (ADR-0064),
    and their strength and novelty are to be re-checked.
  \item \textbf{Choosing the next card from mastery.} Use the mastery estimate to choose what the learner reads
    next and what to review, and evaluate this by mastery rather than engagement (an extension of OBER
    \cite{SRC-134}).
  \item \textbf{Active learning.} Choose which attempts to send to a teacher so that the model improves fastest
    per hour of the teacher's time; this thesis's learning curve is the baseline.
  \item \textbf{Weak labels and continual adaptation.} Use ``the bot was wrong'' disputes and self-corrections as
    weak labels and update the model safely as data accumulate.
  \item \textbf{The effect of feedback.} A controlled experiment: full explanation versus correct/incorrect, with
    skill mastery as the outcome.
  \item \textbf{Transfer and tajweed.} Transfer the dataset, the model and the skill map to speakers of Kazakh,
    Uzbek and Tatar and to the rules of \emph{ikhfa}, \emph{idgham}, \emph{madd} and \emph{qalqala}.
  \item \textbf{An open benchmark.} Release the dataset and the test set as a benchmark of reading by adult
    non-Arab beginners, which does not yet exist.
\end{enumerate}

\section{Closing remark}

Maksudi wrote a book for learners who had a teacher. This thesis does not try to remove the teacher; it gives the
teacher and the learner a careful listener that learns from the teacher to hear better and can say honestly what
the learner has already mastered.
